% GENERATED FILE. Edit canonical lessons, then run npm run generate:books.
% canonical-source-hash: fnv1a64:7cfd0fe6b60fcc99
% canonical-lessons: GU-C165-umervu, GU-C165-redvu, GU-C165-halavvu, GU-C165-bhelavvu, GU-C165-ukalvu

\chapter{Adding, pouring and stirring}
\label{ch:gu-adding-pouring-and-stirring}
\hlchaptermodality{165}

\begin{chapteropening}
\textbf{By the end of this chapter:} \emph{I can describe adding, pouring, stirring and boiling.}

Complete the last lesson of chapter 165: I can describe adding, pouring, stirring and boiling.
\end{chapteropening}

\section[\texorpdfstring{\gu{ઉમેરવું} (umervũ)}{ઉમેરવું (umervũ)}]{\gu{ઉમેરવું} (umervũ) — to add}

\label{lesson:GU-C165-umervu}



\begin{quote}
Before the new one: say the Gujarati for mustard seed, then the Gujarati for cardamom.
\end{quote}

\subsection*{You'll want to know: \gu{ઉમેરવું}}
\textbf{\gu{ઉમેરવું}} — \emph{umervũ} — \textquotedblleft{}to add\textquotedblright{}.

\begin{quote}
\textbf{\gu{તે} \gu{ચામાં} \gu{ખાંડ} \gu{ઉમેરે} \gu{છે}.} — \emph{te cāmā̃ khā̃ḍ umere che} — \textquotedblleft{}He adds sugar to his tea.\textquotedblright{}
\end{quote}

\begin{cousinweb}
\textbf{\gu{ઉમેરવું}} (\emph{umervũ}) is a tadbhava. It comes from \emph{*ummēr}, \textquotedblleft{}join\textquotedblright{}, a reconstructed form, and CDIAL glosses the Gujarati verb \textquotedblleft{}to add, augment\textquotedblright{} (CDIAL 2347).
\end{cousinweb}

\begin{grammarlens}[title={What you've built}]
One more word for everyday life at A2.
\end{grammarlens}

\subsection*{Your turn}
\begin{itemize}
\raggedright
  \item \emph{Say it:} \emph{umervũ}
  \item \emph{Say it:} \emph{umervũ}, once more
  \item \emph{Say it:} \emph{te cāmā̃ khā̃ḍ umere che} — He adds sugar to his tea.
  \item \emph{Recall:} say the Gujarati for mustard seed, then the Gujarati for cardamom, then say \emph{umervũ} again
\end{itemize}

\begin{tcolorbox}[breakable,colback=teal!4,colframe=teal!35!black,title={Before you move on}]
What does \gu{ઉમેરવું} mean? (\textquotedblleft{}To add\textquotedblright{}.) Say it once more.
\end{tcolorbox}
\section[\texorpdfstring{\gu{રેડવું} (reḍvũ)}{રેડવું (reḍvũ)}]{\gu{રેડવું} (reḍvũ) — to pour}

\label{lesson:GU-C165-redvu}



\begin{quote}
Before the new one: say the Gujarati for cardamom, then the Gujarati for to add.
\end{quote}

\subsection*{You'll want to know: \gu{રેડવું}}
\textbf{\gu{રેડવું}} — \emph{reḍvũ} — \textquotedblleft{}to pour\textquotedblright{}.

\begin{quote}
\textbf{\gu{હું} \gu{પ્યાલામાં} \gu{ચા} \gu{રેડું} \gu{છું}.} — \emph{hũ pyālāmā̃ cā reḍũ chũ} — \textquotedblleft{}I pour tea into the cup.\textquotedblright{}
\end{quote}

\begin{cousinweb}
\textbf{\gu{રેડવું}} (\emph{reḍvũ}) is inherited, but its source is not settled. CDIAL lists it, glossed \textquotedblleft{}to pour\textquotedblright{}, under Sanskrit \emph{raya}, \textquotedblleft{}stream\textquotedblright{}, and again under \emph{recayati}, \textquotedblleft{}empties\textquotedblright{}, and says it may come from either (CDIAL 10638, 10811).
\end{cousinweb}

\begin{grammarlens}[title={What you've built}]
One more word for everyday life at A2.
\end{grammarlens}

\subsection*{Your turn}
\begin{itemize}
\raggedright
  \item \emph{Say it:} \emph{reḍvũ}
  \item \emph{Say it:} \emph{reḍvũ}, once more
  \item \emph{Say it:} \emph{hũ pyālāmā̃ cā reḍũ chũ} — I pour tea into the cup.
  \item \emph{Recall:} say the Gujarati for cardamom, then the Gujarati for to add, then say \emph{reḍvũ} again
\end{itemize}

\begin{tcolorbox}[breakable,colback=teal!4,colframe=teal!35!black,title={Before you move on}]
What does \gu{રેડવું} mean? (\textquotedblleft{}To pour\textquotedblright{}.) Say it once more.
\end{tcolorbox}
\section[\texorpdfstring{\gu{હલાવવું} (halāvvũ)}{હલાવવું (halāvvũ)}]{\gu{હલાવવું} (halāvvũ) — to stir, to shake}

\label{lesson:GU-C165-halavvu}



\begin{quote}
Before the new one: say the Gujarati for to add, then the Gujarati for to pour.
\end{quote}

\subsection*{You'll want to know: \gu{હલાવવું}}
\textbf{\gu{હલાવવું}} — \emph{halāvvũ} — \textquotedblleft{}to stir, to shake\textquotedblright{}.

\begin{quote}
\textbf{\gu{તે} \gu{ચમચીથી} \gu{દાળ} \gu{હલાવે} \gu{છે}.} — \emph{te camcīthī dāḷ halāve che} — \textquotedblleft{}She stirs the dal with a spoon.\textquotedblright{}
\end{quote}

\begin{cousinweb}
\textbf{\gu{હલાવવું}} (\emph{halāvvũ}) is a tadbhava, a causative. CDIAL lists it, glossed \textquotedblleft{}to shake\textquotedblright{}, under \emph{*hallati}, \textquotedblleft{}moves, shakes\textquotedblright{}, a reconstructed form (CDIAL 14018).
\end{cousinweb}

\begin{grammarlens}[title={What you've built}]
One more word for everyday life at A2.
\end{grammarlens}

\subsection*{Your turn}
\begin{itemize}
\raggedright
  \item \emph{Say it:} \emph{halāvvũ}
  \item \emph{Say it:} \emph{halāvvũ}, once more
  \item \emph{Say it:} \emph{te camcīthī dāḷ halāve che} — She stirs the dal with a spoon.
  \item \emph{Recall:} say the Gujarati for to add, then the Gujarati for to pour, then say \emph{halāvvũ} again
\end{itemize}

\begin{tcolorbox}[breakable,colback=teal!4,colframe=teal!35!black,title={Before you move on}]
What does \gu{હલાવવું} mean? (\textquotedblleft{}To stir, to shake\textquotedblright{}.) Say it once more.
\end{tcolorbox}
\section[\texorpdfstring{\gu{ભેળવવું} (bheḷavvũ)}{ભેળવવું (bheḷavvũ)}]{\gu{ભેળવવું} (bheḷavvũ) — to mix}

\label{lesson:GU-C165-bhelavvu}



\begin{quote}
Before the new one: say the Gujarati for to pour, then the Gujarati for to stir.
\end{quote}

\subsection*{You'll want to know: \gu{ભેળવવું}}
\textbf{\gu{ભેળવવું}} — \emph{bheḷavvũ} — \textquotedblleft{}to mix\textquotedblright{}.

\begin{quote}
\textbf{\gu{તે} \gu{લોટમાં} \gu{પાણી} \gu{ભેળવે} \gu{છે}.} — \emph{te loṭmā̃ pāṇī bheḷave che} — \textquotedblleft{}He mixes water into the flour.\textquotedblright{}
\end{quote}

\begin{cousinweb}
\textbf{\gu{ભેળવવું}} (\emph{bheḷavvũ}) is a tadbhava. CDIAL lists it, glossed \textquotedblleft{}to mix\textquotedblright{}, under \emph{*bhiṭ}, \emph{*bhēl}, \textquotedblleft{}meet, throng\textquotedblright{}, a group of reconstructed forms (CDIAL 9490).
\end{cousinweb}

\begin{grammarlens}[title={What you've built}]
One more word for everyday life at A2.
\end{grammarlens}

\subsection*{Your turn}
\begin{itemize}
\raggedright
  \item \emph{Say it:} \emph{bheḷavvũ}
  \item \emph{Say it:} \emph{bheḷavvũ}, once more
  \item \emph{Say it:} \emph{te loṭmā̃ pāṇī bheḷave che} — He mixes water into the flour.
  \item \emph{Recall:} say the Gujarati for to pour, then the Gujarati for to stir, then say \emph{bheḷavvũ} again
\end{itemize}

\begin{tcolorbox}[breakable,colback=teal!4,colframe=teal!35!black,title={Before you move on}]
What does \gu{ભેળવવું} mean? (\textquotedblleft{}To mix\textquotedblright{}.) Say it once more.
\end{tcolorbox}
\section[\texorpdfstring{\gu{ઉકાળવું} (ukāḷvũ)}{ઉકાળવું (ukāḷvũ)}]{\gu{ઉકાળવું} (ukāḷvũ) — to boil}

\label{lesson:GU-C165-ukalvu}



\begin{quote}
Before the new one: say the Gujarati for to stir, then the Gujarati for to mix.
\end{quote}

\subsection*{You'll want to know: \gu{ઉકાળવું}}
\textbf{\gu{ઉકાળવું}} — \emph{ukāḷvũ} — \textquotedblleft{}to boil\textquotedblright{}.

\begin{quote}
\textbf{\gu{માતા} \gu{સવારે} \gu{દૂધ} \gu{ઉકાળે} \gu{છે}.} — \emph{mātā savāre dūdh ukāḷe che} — \textquotedblleft{}Mother boils the milk in the morning.\textquotedblright{}
\end{quote}

\begin{cousinweb}
\textbf{\gu{ઉકાળવું}} (\emph{ukāḷvũ}) is a tadbhava. CDIAL lists it, glossed \textquotedblleft{}to cause to boil\textquotedblright{}, under Sanskrit \emph{utkālayati}, \textquotedblleft{}drives out\textquotedblright{}, in the entry for \emph{*utkalati}, \textquotedblleft{}goes out or up\textquotedblright{}, a reconstructed form that gives the intransitive \emph{ukaḷvũ}, \textquotedblleft{}to boil\textquotedblright{} (CDIAL 1716).
\end{cousinweb}

\begin{grammarlens}[title={What you've built}]
One more word for everyday life at A2.
\end{grammarlens}

\subsection*{Your turn}
\begin{itemize}
\raggedright
  \item \emph{Say it:} \emph{ukāḷvũ}
  \item \emph{Say it:} \emph{ukāḷvũ}, once more
  \item \emph{Say it:} \emph{mātā savāre dūdh ukāḷe che} — Mother boils the milk in the morning.
  \item \emph{Recall:} say the Gujarati for to stir, then the Gujarati for to mix, then say \emph{ukāḷvũ} again
\end{itemize}

\begin{tcolorbox}[breakable,colback=teal!4,colframe=teal!35!black,title={Before you move on}]
What does \gu{ઉકાળવું} mean? (\textquotedblleft{}To boil\textquotedblright{}.) Say it once more.
\end{tcolorbox}
